% GENERATED FILE. Edit canonical lessons, then run npm run generate:books.
% canonical-source-hash: fnv1a64:b44d2e57432cac01
% canonical-lessons: FA-C285-fizik, FA-C285-shimi, FA-C285-adabiyat, FA-C285-engelisi, FA-C285-hendese

\chapter{Physics, Chemistry, Literature, English, Geometry}
\label{ch:fa-physics-chemistry-literature-english-geometry}
\hlchaptermodality{285}

\begin{chapteropening}
\textbf{By the end of this chapter:} \emph{I can say physics, chemistry, literature, English and geometry.}

Complete the last lesson of chapter 285: I can say physics, chemistry, literature, English and geometry.
\end{chapteropening}

\section[\texorpdfstring{\fa{فیزیک} (fizik)}{فیزیک (fizik)}]{\fa{فیزیک} (fizik) — physics}

\label{lesson:FA-C285-fizik}



\begin{quote}
Before the new one: say the Persian for a librarian, then the Persian for an accountant.
\end{quote}

\subsection*{You'll want to know: \fa{فیزیک}}
\textbf{\fa{فیزیک}} — \emph{fizik} — \textquotedblleft{}physics\textquotedblright{}.

\begin{cousinweb}
\emph{fizik} is a loan, usually traced to French \emph{physique}, \textquotedblleft{}physics\textquotedblright{}.
\end{cousinweb}

\begin{grammarlens}[title={What you've built}]
One more word for this chapter.
\end{grammarlens}

\subsection*{Your turn}
\begin{itemize}
\raggedright
  \item \emph{Say it:} \emph{fizik}
  \item \emph{Say it:} \emph{fizik}, once more
  \item \emph{Say it:} say \emph{fizik}
  \item \emph{Recall:} say the Persian for a librarian, then the Persian for an accountant, then say \emph{fizik} again
\end{itemize}

\begin{tcolorbox}[breakable,colback=teal!4,colframe=teal!35!black,title={Before you move on}]
What does \fa{فیزیک} mean? (\textquotedblleft{}Physics\textquotedblright{}.) Say it once more.
\end{tcolorbox}
\section[\texorpdfstring{\fa{شیمی} (shimi)}{شیمی (shimi)}]{\fa{شیمی} (shimi) — chemistry}

\label{lesson:FA-C285-shimi}



\begin{quote}
Before the new one: say the Persian for an accountant, then the Persian for physics.
\end{quote}

\subsection*{You'll want to know: \fa{شیمی}}
\textbf{\fa{شیمی}} — \emph{shimi} — \textquotedblleft{}chemistry\textquotedblright{}.

\begin{cousinweb}
\emph{shimi} is borrowed from French \emph{chimie}, \textquotedblleft{}chemistry\textquotedblright{} (EtymDB).
\end{cousinweb}

\begin{grammarlens}[title={What you've built}]
One more word for this chapter.
\end{grammarlens}

\subsection*{Your turn}
\begin{itemize}
\raggedright
  \item \emph{Say it:} \emph{shimi}
  \item \emph{Say it:} \emph{shimi}, once more
  \item \emph{Say it:} say \emph{shimi}
  \item \emph{Recall:} say the Persian for an accountant, then the Persian for physics, then say \emph{shimi} again
\end{itemize}

\begin{tcolorbox}[breakable,colback=teal!4,colframe=teal!35!black,title={Before you move on}]
What does \fa{شیمی} mean? (\textquotedblleft{}Chemistry\textquotedblright{}.) Say it once more.
\end{tcolorbox}
\section[\texorpdfstring{\fa{ادبیات} (adabiyât)}{ادبیات (adabiyât)}]{\fa{ادبیات} (adabiyât) — literature}

\label{lesson:FA-C285-adabiyat}



\begin{quote}
Before the new one: say the Persian for physics, then the Persian for chemistry.
\end{quote}

\subsection*{You'll want to know: \fa{ادبیات}}
\textbf{\fa{ادبیات}} — \emph{adabiyât} — \textquotedblleft{}literature\textquotedblright{}.

\begin{cousinweb}
\emph{adabiyât} is borrowed from Arabic \emph{adabiyyāt}, \textquotedblleft{}literature, letters\textquotedblright{}. Its Arabic root is \emph{ʾ-d-b} (Lane).
\end{cousinweb}

\begin{grammarlens}[title={What you've built}]
One more word for this chapter.
\end{grammarlens}

\subsection*{Your turn}
\begin{itemize}
\raggedright
  \item \emph{Say it:} \emph{adabiyât}
  \item \emph{Say it:} \emph{adabiyât}, once more
  \item \emph{Say it:} say \emph{adabiyât}
  \item \emph{Recall:} say the Persian for physics, then the Persian for chemistry, then say \emph{adabiyât} again
\end{itemize}

\begin{tcolorbox}[breakable,colback=teal!4,colframe=teal!35!black,title={Before you move on}]
What does \fa{ادبیات} mean? (\textquotedblleft{}Literature\textquotedblright{}.) Say it once more.
\end{tcolorbox}
\section[\texorpdfstring{\fa{انگلیسی} (engelisi)}{انگلیسی (engelisi)}]{\fa{انگلیسی} (engelisi) — English (the language)}

\label{lesson:FA-C285-engelisi}



\begin{quote}
Before the new one: say the Persian for chemistry, then the Persian for literature.
\end{quote}

\subsection*{You'll want to know: \fa{انگلیسی}}
\textbf{\fa{انگلیسی}} — \emph{engelisi} — \textquotedblleft{}English (the language)\textquotedblright{}.

\begin{cousinweb}
\emph{engelisi} is \emph{Engelis}, \textquotedblleft{}England, the English\textquotedblright{}, with the ending \emph{-i}. \emph{Engelis} is a loan from European languages, whose name for the English appears as French \emph{anglais} and Italian \emph{inglese}; which one passed it to Persian is not certain.
\end{cousinweb}

\begin{grammarlens}[title={What you've built}]
One more word for this chapter.
\end{grammarlens}

\subsection*{Your turn}
\begin{itemize}
\raggedright
  \item \emph{Say it:} \emph{engelisi}
  \item \emph{Say it:} \emph{engelisi}, once more
  \item \emph{Say it:} say \emph{engelisi}
  \item \emph{Recall:} say the Persian for chemistry, then the Persian for literature, then say \emph{engelisi} again
\end{itemize}

\begin{tcolorbox}[breakable,colback=teal!4,colframe=teal!35!black,title={Before you move on}]
What does \fa{انگلیسی} mean? (\textquotedblleft{}English (the language)\textquotedblright{}.) Say it once more.
\end{tcolorbox}
\section[\texorpdfstring{\fa{هندسه} (hendese)}{هندسه (hendese)}]{\fa{هندسه} (hendese) — geometry}

\label{lesson:FA-C285-hendese}



\begin{quote}
Before the new one: say the Persian for literature, then the Persian for English.
\end{quote}

\subsection*{You'll want to know: \fa{هندسه}}
\textbf{\fa{هندسه}} — \emph{hendese} — \textquotedblleft{}geometry\textquotedblright{}.

\begin{cousinweb}
\emph{hendese} is borrowed from Arabic \emph{handasa}, \textquotedblleft{}geometry\textquotedblright{}, which EtymDB traces back to Persian \emph{andâze}, \textquotedblleft{}measure\textquotedblright{}: a Persian word home again from Arabic.
\end{cousinweb}

\begin{grammarlens}[title={What you've built}]
One more word for this chapter.
\end{grammarlens}

\subsection*{Your turn}
\begin{itemize}
\raggedright
  \item \emph{Say it:} \emph{hendese}
  \item \emph{Say it:} \emph{hendese}, once more
  \item \emph{Say it:} say \emph{hendese}
  \item \emph{Recall:} say the Persian for literature, then the Persian for English, then say \emph{hendese} again
\end{itemize}

\begin{tcolorbox}[breakable,colback=teal!4,colframe=teal!35!black,title={Before you move on}]
What does \fa{هندسه} mean? (\textquotedblleft{}Geometry\textquotedblright{}.) Say it once more.
\end{tcolorbox}
